% !TEX program = pdflatex
% Анализ, верификация и сопоставление реализации непараметрической знаковой
% карты Шухарта (SN) со средствами языка R.
% Первоисточники: Amin, Reynolds & Bakir (1995); Gibbons & Chakraborti (2010);
% Nascimento et al. (2026). Сборка: pdflatex (T2A + babel russian).
\documentclass[11pt,a4paper]{article}

\usepackage[T2A]{fontenc}
\usepackage[utf8]{inputenc}
\usepackage[russian]{babel}
\usepackage[margin=2.4cm]{geometry}
\usepackage{amsmath,amssymb}
\usepackage{pgfplots}
\pgfplotsset{compat=1.17}
\usepackage{booktabs}
\usepackage{array}
\usepackage{graphicx}
\usepackage{xcolor}
\usepackage{listings}
\usepackage{caption}
\usepackage{enumitem}
\usepackage[hidelinks]{hyperref}

\definecolor{accent}{RGB}{46,134,171}   % ACCENT  #2E86AB
\definecolor{alert}{RGB}{178,34,34}     % ALERT   #B22222
\definecolor{third}{RGB}{60,110,71}     % PALETTE #3C6E47
\definecolor{softgray}{RGB}{128,128,128} % CENTER #808080

\lstset{
  basicstyle=\ttfamily\small,
  commentstyle=\itshape\color{softgray},
  keywordstyle=\bfseries,
  breaklines=true,
  frame=single,
  rulecolor=\color{softgray!50},
  columns=fullflexible,
  keepspaces=true,
  showstringspaces=false,
  aboveskip=6pt,belowskip=6pt,
}

\captionsetup{font=small,labelfont=bf}

% Небольшой допуск на растяжение строк: предотвращает редкие переполнения
% из-за длинных неразрывных имён методов в \texttt{}.
\setlength{\emergencystretch}{2.5em}

\newcommand{\SN}{\mathrm{SN}}
\DeclareMathOperator*{\argmin}{arg\,min}

\pgfplotsset{
  myaxis/.style={
    width=0.92\textwidth, height=0.42\textwidth,
    grid=both, grid style={softgray!25},
    axis line style={softgray}, tick align=outside,
    label style={font=\small}, tick label style={font=\small},
    legend style={font=\footnotesize, draw=softgray!50, fill=white},
    every axis plot/.append style={line width=1pt},
  },
}

% --- Координаты графиков (см. experiments/sn/gen_figdata.py) ---
\def\verACons{(1,1.000000) (2,2.000000) (3,3.000000) (4,4.000000) (5,5.000000) (6,6.000000) (7,7.000000) (8,8.000000) (9,9.000000) (10,10.000000) (11,11.000000) (12,12.000000) (13,13.000000) (14,12.000000) (15,13.000000) (16,14.000000) (17,13.000000) (18,14.000000) (19,15.000000) (20,14.000000)}
\def\verANp{(1,1.000000) (2,2.000000) (3,3.000000) (4,4.000000) (5,5.000000) (6,6.000000) (7,7.000000) (8,8.000000) (9,9.000000) (10,10.000000) (11,11.000000) (12,12.000000) (13,11.000000) (14,12.000000) (15,13.000000) (16,12.000000) (17,13.000000) (18,14.000000) (19,13.000000) (20,14.000000)}
\def\verANear{(1,1.000000) (2,1.000000) (3,2.000000) (4,3.000000) (5,4.000000) (6,5.000000) (7,6.000000) (8,5.000000) (9,6.000000) (10,7.000000) (11,6.000000) (12,7.000000) (13,8.000000) (14,7.000000) (15,8.000000) (16,7.000000) (17,8.000000) (18,9.000000) (19,8.000000) (20,9.000000)}
\def\verOCBook{(0.0,0.000519) (0.25,0.003207) (0.5,0.022228) (0.75,0.093159) (1.0,0.253224) (1.25,0.484208) (1.5,0.709673) (1.75,0.867017) (2.0,0.949718)}
\def\verOCNp{(0.0,0.004181) (0.25,0.018020) (0.5,0.087478) (0.75,0.261912) (1.0,0.522208) (1.25,0.764768) (1.5,0.913109) (1.75,0.975662) (2.0,0.994720)}
% --- режимы выбора a ---
\def\verModeCons{(1,1.000000) (2,2.000000) (3,3.000000) (4,4.000000) (5,5.000000) (6,6.000000) (7,7.000000) (8,8.000000) (9,9.000000) (10,10.000000) (11,11.000000) (12,12.000000) (13,13.000000) (14,12.000000) (15,13.000000) (16,14.000000) (17,13.000000) (18,14.000000) (19,15.000000) (20,14.000000) (21,15.000000) (22,16.000000) (23,15.000000) (24,16.000000) (25,17.000000) (26,16.000000) (27,17.000000) (28,18.000000) (29,17.000000) (30,18.000000)}
\def\verModeSym{(1,1.000000) (2,2.000000) (3,3.000000) (4,4.000000) (5,5.000000) (6,6.000000) (7,7.000000) (8,8.000000) (9,9.000000) (10,10.000000) (11,11.000000) (12,12.000000) (13,11.000000) (14,12.000000) (15,13.000000) (16,12.000000) (17,13.000000) (18,14.000000) (19,15.000000) (20,14.000000) (21,15.000000) (22,16.000000) (23,15.000000) (24,16.000000) (25,15.000000) (26,16.000000) (27,17.000000) (28,16.000000) (29,17.000000) (30,18.000000)}
\def\verModeNear{(1,1.000000) (2,1.000000) (3,2.000000) (4,3.000000) (5,4.000000) (6,5.000000) (7,6.000000) (8,5.000000) (9,6.000000) (10,7.000000) (11,6.000000) (12,7.000000) (13,8.000000) (14,7.000000) (15,8.000000) (16,7.000000) (17,8.000000) (18,9.000000) (19,8.000000) (20,9.000000) (21,10.000000) (22,9.000000) (23,10.000000) (24,9.000000) (25,10.000000) (26,11.000000) (27,10.000000) (28,11.000000) (29,10.000000) (30,11.000000)}
\def\verFARCons{(1,1.000000) (2,0.500000) (3,0.250000) (4,0.125000) (5,0.062500) (6,0.031250) (7,0.015625) (8,0.007812) (9,0.003906) (10,0.001953) (11,0.000977) (12,0.000488) (13,0.000244) (14,0.001831) (15,0.000977) (16,0.000519) (17,0.002350) (18,0.001312) (19,0.000729) (20,0.002577) (21,0.001490) (22,0.000855) (23,0.002599) (24,0.001544) (25,0.000911) (26,0.002494) (27,0.001514) (28,0.000912) (29,0.002316) (30,0.001431)}
\def\verFARSym{(1,1.000000) (2,0.500000) (3,0.250000) (4,0.125000) (5,0.062500) (6,0.031250) (7,0.015625) (8,0.007812) (9,0.003906) (10,0.001953) (11,0.000977) (12,0.000488) (13,0.003418) (14,0.001831) (15,0.000977) (16,0.004181) (17,0.002350) (18,0.001312) (19,0.000729) (20,0.002577) (21,0.001490) (22,0.000855) (23,0.002599) (24,0.001544) (25,0.004077) (26,0.002494) (27,0.001514) (28,0.003719) (29,0.002316) (30,0.001431)}
\def\verFARNear{(1,1.000000) (2,0.500000) (3,0.250000) (4,0.125000) (5,0.062500) (6,0.031250) (7,0.015625) (8,0.070312) (9,0.039062) (10,0.021484) (11,0.065430) (12,0.038574) (13,0.022461) (14,0.057373) (15,0.035156) (16,0.076813) (17,0.049042) (18,0.030884) (19,0.063568) (20,0.041389) (21,0.026604) (22,0.052479) (23,0.034690) (24,0.063915) (25,0.043285) (26,0.028959) (27,0.052239) (28,0.035698) (29,0.061428) (30,0.042774)}
\def\verModeTable{%
1 & 1 & 1.000000 & 1 & 1.000000 & 1 & 1.000000 \\
2 & 2 & 0.500000 & 2 & 0.500000 & 1 & 0.500000 \\
3 & 3 & 0.250000 & 3 & 0.250000 & 2 & 0.250000 \\
4 & 4 & 0.125000 & 4 & 0.125000 & 3 & 0.125000 \\
5 & 5 & 0.062500 & 5 & 0.062500 & 4 & 0.062500 \\
6 & 6 & 0.031250 & 6 & 0.031250 & 5 & 0.031250 \\
7 & 7 & 0.015625 & 7 & 0.015625 & 6 & 0.015625 \\
8 & 8 & 0.007812 & 8 & 0.007812 & 5 & 0.070312 \\
9 & 9 & 0.003906 & 9 & 0.003906 & 6 & 0.039062 \\
10 & 10 & 0.001953 & 10 & 0.001953 & 7 & 0.021484 \\
11 & 11 & 0.000977 & 11 & 0.000977 & 6 & 0.065430 \\
12 & 12 & 0.000488 & 12 & 0.000488 & 7 & 0.038574 \\
13 & 13 & 0.000244 & 11 & 0.003418 & 8 & 0.022461 \\
14 & 12 & 0.001831 & 12 & 0.001831 & 7 & 0.057373 \\
15 & 13 & 0.000977 & 13 & 0.000977 & 8 & 0.035156 \\
16 & 14 & 0.000519 & 12 & 0.004181 & 7 & 0.076813 \\
17 & 13 & 0.002350 & 13 & 0.002350 & 8 & 0.049042 \\
18 & 14 & 0.001312 & 14 & 0.001312 & 9 & 0.030884 \\
19 & 15 & 0.000729 & 15 & 0.000729 & 8 & 0.063568 \\
20 & 14 & 0.002577 & 14 & 0.002577 & 9 & 0.041389 \\}
\def\verNearAbs14{(1,0.790527) (2,0.790527) (3,0.423950) (4,0.423950) (5,0.179565) (6,0.179565) (7,0.057373) (8,0.057373) (9,0.012939) (10,0.012939) (11,0.001831) (12,0.001831) (13,0.000122) (14,0.000122)}
\def\verNearMin14{(7,0.057373) (8,0.057373)}

\title{\textbf{Анализ, верификация и сопоставление реализации}\\[2pt]
       \textbf{непараметрической знаковой карты Шухарта (SN)}}
\author{Проект \texttt{nonparametric-charts}}
\date{\today}

\begin{document}
\maketitle

% ================================================================== %
\section*{Сокращения}
\addcontentsline{toc}{section}{Сокращения}
% ================================================================== %

\begin{description}[leftmargin=2.6cm,itemsep=1pt]
  \item[SN] (sign statistic) --- знаковая статистика;
  \item[CL] (center line) --- центральная линия;
  \item[UCL, LCL] (upper/lower control limit) --- верхняя и нижняя контрольные границы;
  \item[FAR] (false alarm rate) --- вероятность ложной тревоги на одной подгруппе;
  \item[ARL] (average run length) --- средняя длина серии;
  \item[MRL] (median run length) --- медианная длина серии;
  \item[SD] (standard deviation) --- среднеквадратичное отклонение;
  \item[IC, OOC] (in-control, out-of-control) --- управляемое и неуправляемое состояние процесса;
  \item[$\mu_0$] --- целевое (референтное) значение центра процесса;
  \item[n, m] --- размер подгруппы и число подгрупп.
\end{description}

\tableofcontents
\bigskip

% ================================================================== %
\section{Математические основы реализации}
% ================================================================== %

\subsection{Знаковая статистика}

Для $i$-й подгруппы объёма $n$ со значениями $X_{i1},\dots,X_{in}$ и целевым
значением $\mu_0$ знаковая статистика определяется как
\begin{equation}
  \SN_i = n_i^{+} - n_i^{-}
        = \sum_{j=1}^{n}\phi(X_{ij}-\mu_0)
        = 2T_i - n,
  \qquad
  \phi(t) = \begin{cases} 1, & t>0,\\ 0, & t=0,\\ -1, & t<0,\end{cases}
  \label{eq:sn}
\end{equation}
где $n_i^{+}$ и $n_i^{-}$ --- числа наблюдений, больших и меньших $\mu_0$, а
$T_i = n_i^{+}$ --- классическая знаковая статистика знакового критерия
\cite{Amin1995, Gibbons2010}. В управляемом состоянии
$T_i \sim \mathrm{Bin}(n,\tfrac12)$, вследствие чего $\SN_i$ принимает
целочисленные значения той же чётности, что и $n$, и
$\mathbb{E}[\SN_i]=0$.

Границы полагаются симметричными:
\begin{equation}
  \mathrm{UCL} = a, \qquad \mathrm{CL} = 0, \qquad \mathrm{LCL} = -a,
  \qquad 0<a\le n,
  \label{eq:limits}
\end{equation}
а процесс признаётся неуправляемым при
\begin{equation}
  |\SN_i| \ge a.
  \label{eq:signal}
\end{equation}
Такая запись соответствует постановке задачи
\cite{Amin1995, Chakraborti2019}.

\subsection{Выбор константы $a$ и дискретность}

Постоянная $a$ выбирается как наименьшее целое $a$ ($0<a\le n$) той же
чётности, что и $n$, для которого
\begin{equation}
  P_{\mu_0|\mathrm{IC}}\bigl(\SN_i \ge a\bigr) \le \frac{\mathrm{FAR}_0}{2},
  \qquad \mathrm{FAR}_0 = 0.0027,
  \label{eq:rule}
\end{equation}
что обеспечивает $\mathrm{FAR} = 2P(\SN_i\ge a) \le \mathrm{FAR}_0$ и
$\mathrm{ARL} = 1/\mathrm{FAR}$. Номинальную величину можно задавать
эквивалентно через $\mathrm{ARL}_0=1/\mathrm{FAR}_0\approx370$ значение $\mathrm{FAR}_0=0.0027$ приняты в литературе 
по знаковым картам \cite{Amin1995, Chakraborti2019}.

Поскольку распределение $\SN_i$ дискретно, равенство
\eqref{eq:rule} обычно не достигается. Реализация применяет
\emph{консервативное} округление в сторону ужесточения: выбирается
наименьшее достижимое значение $a$, при котором неравенство выполняется;
если же при малом $n$ таких значений не существует, полагается $a=n$
(наиболее широкие границы и наименьшая достижимая FAR). Дополнительно реализован режим
\texttt{nearest} ($a=\argmin_a |P(|\SN_i|\ge a)-\alpha|$, $\alpha=0.05$),
имеющий демонстрационное назначение: он не учитывает чётность и потому может
возвращать значения, недостижимые для $\SN_i$. Подробнее методы выбора
константы $a$ описаны в подразделе~\ref{subsec:modes}.

Характеристики длины серии при отсутствии ложных тревог распределены
геометрически, откуда
$\mathrm{ARL}=1/\mathrm{FAR}$ и
$\mathrm{MRL}=\lceil \ln 0.5 / \ln(1-\mathrm{FAR})\rceil$
\cite{Chakraborti2019}.

\subsection{Методы выбора константы $a$}
\label{subsec:modes}

Библиотека реализует три метода выбора $a$; переключение выполняется
параметром \texttt{approach}. Основным является точный метод,
реализующий правило~\eqref{eq:rule}; режим \texttt{normal} --- асимптотическая
(нормальная) аппроксимация, метод «ближайший» --- демонстрационный.

\paragraph{Точный метод.}
Константа выбирается как наименьшее достижимое значение
\begin{equation}
  a_{\mathrm{exact}}=\min\{\,a:\ 1\le a\le n,\ a\equiv n\ (\mathrm{mod}\ 2),\
  P(\SN_i\ge a)\le \mathrm{FAR}_0/2\,\},
  \label{eq:acons}
\end{equation}
а при отсутствии такого значения полагается $a=n$ (наиболее широкие границы
и наименьшая достижимая FAR). Выбор в сторону ужесточения гарантирует
$\mathrm{FAR}=2P(\SN_i\ge a)\le \mathrm{FAR}_0$. Требование $a\equiv n\
(\mathrm{mod}\ 2)$ следует из дискретности распределения $\SN_i$: прочие
значения $a$ недостижимы и дали бы ту же самую границу
\cite{Amin1995, Gibbons2010}. Односторонняя вероятность, что статистика окажется не меньше верхней границы или «хвост» вычисляется по следующей формуле:
\begin{equation}
  P(\SN_i\ge a)
  =\sum_{m=\lceil (n+a)/2\rceil}^{n}\binom{n}{m}\Bigl(\frac12\Bigr)^{\!n}.
  \label{eq:sntail}
\end{equation}

\paragraph{Нормальная аппроксимация.}
Граница берётся по нормальному приближению распределения $\SN_i$ в IC,
$SD(\SN_i)=\sqrt n$:
\begin{equation}
  a_{\mathrm{norm}}=\min\{\,a:\ 1\le a\le n,\ a\equiv n\ (\mathrm{mod}\ 2),\
  a\ge z_{1-\mathrm{FAR}_0/2}\,\sqrt n\,\},
  \label{eq:asym}
\end{equation}
где $z_{1-\mathrm{FAR}_0/2}\approx2.9999$ --- квантиль стандартного
нормального закона. Риск делится симметрично: по $z_{1-\mathrm{FAR}_0/2}\sqrt n$
на каждый «хвост». Это асимптотический по природе метод (ср.\ режим
\texttt{asymptotic} в \texttt{npSPC}), поэтому достигнутый FAR не обязан
быть не больше $\mathrm{FAR}_0$; при малых $n$, где порог превышает $n$,
полагается $a=n$.

\paragraph{Метод «ближайший».}
Константа выбирается как ближайшая к заданному уровню значимости
$\alpha=0.05$:
\begin{equation}
  a_{\mathrm{near}}=\argmin_{1\le a\le n}\bigl|P(|\SN_i|\ge a)-\alpha\bigr|,
  \qquad P(|\SN_i|\ge a)=2\,P(\SN_i\ge a).
  \label{eq:anear}
\end{equation}
Перебор идёт по всем целым $a$ без требования чётности, поэтому
$a_{\mathrm{near}}$ может не быть достижимым значением $\SN_i$; тогда
фактическая граница равна соседнему достижимому значению (для симметричных
границ $\pm a$ это эквивалентно). Метод имеет демонстрационное назначение: к
статистическому управлению процессами он не применим, поскольку
$\alpha=0.05$ отвечает $\mathrm{FAR}\approx0.05$ и $\mathrm{ARL}\approx20$,
что далеко от принятого $\mathrm{FAR}_0=0.0027$.

Сопоставление методов при $n=1,\dots,20$ приведено в
таблице~\ref{tab:modes}. Точный метод и нормальная аппроксимация
совпадают при малых $n$ и расходятся там, где дискретность делает
точную границу заметно шире асимптотической; метод «ближайший»
даёт самые узкие границы.

\begin{table}[htbp]\centering
\caption{Методы выбора константы $a$, $n=1..20$
($\mathrm{FAR}_0=0.0027$, $\alpha=0.05$).}
\label{tab:modes}
\small
\begin{tabular}{@{}r r r r r r r@{}}
\toprule
$n$ & $a_{\mathrm{exact}}$ & $\mathrm{FAR}_{\mathrm{exact}}$
& $a_{\mathrm{norm}}$ & $\mathrm{FAR}_{\mathrm{norm}}$
& $a_{\mathrm{near}}$ & $\mathrm{FAR}_{\mathrm{near}}$ \\
\midrule
\verModeTable
\bottomrule
\end{tabular}
\end{table}

\begin{figure}[htbp]\centering
\begin{tikzpicture}
\begin{axis}[myaxis,
  xlabel={$n$}, ylabel={константа $a$},
  xmin=0, xmax=31, ymin=0, ymax=19,
  legend pos=north west,
  title={Константа $a(n)$ по методам}]
\addplot+[mark=*,mark size=1.4pt,color=accent] coordinates {\verModeCons};
\addlegendentry{точный}
\addplot+[mark=square*,mark size=1.4pt,color=third] coordinates {\verModeSym};
\addlegendentry{нормальная аппроксимация}
\addplot+[mark=triangle,mark size=1.6pt,color=alert] coordinates {\verModeNear};
\addlegendentry{ближайший ($\alpha=0.05$)}
\end{axis}
\end{tikzpicture}
\caption{Константа $a(n)$ для трёх методов. Точный метод целится в
$\mathrm{FAR}_0/2=0.00135$, нормальная аппроксимация --- в
$z_{1-\mathrm{FAR}_0/2}\sqrt n$, метод «ближайший» --- в $\alpha=0.05$.}
\label{fig:modesa}
\end{figure}

\begin{figure}[htbp]\centering
\begin{tikzpicture}
\begin{axis}[myaxis,
  xlabel={$n$}, ylabel={достигнутый $\mathrm{FAR}$},
  xmin=0, xmax=31, ymode=log, log basis y=10,
  ymin=1e-4, ymax=1.2,
  legend pos=north east,
  title={Достигнутый $\mathrm{FAR}(n)$ по методам}]
\addplot+[mark=*,mark size=1.4pt,color=accent] coordinates {\verFARCons};
\addlegendentry{точный}
\addplot+[mark=square*,mark size=1.4pt,color=third] coordinates {\verFARSym};
\addlegendentry{нормальная аппроксимация}
\addplot+[mark=triangle,mark size=1.6pt,color=alert] coordinates {\verFARNear};
\addlegendentry{ближайший}
\addplot[dashed,black,line width=1pt,domain=0:31] {0.0027};
\addlegendentry{$\mathrm{FAR}_0=0.0027$}
\addplot[dashed,alert,line width=1pt,domain=0:31] {0.05};
\addlegendentry{$\alpha=0.05$}
\end{axis}
\end{tikzpicture}
\caption{Достигнутый двусторонний $\mathrm{FAR}$ при трёх методах. Кривая
точного метода всюду лежит ниже $\mathrm{FAR}_0$ (с «пилообразными»
скачками из-за дискретности); кривая нормальной аппроксимации колеблется
вокруг $\mathrm{FAR}_0$, а кривая метода «ближайший» --- вокруг $\alpha$ и на
порядок выше.}
\label{fig:modesfar}
\end{figure}

\paragraph{Пример $n=14$.}
Возможные значения $\SN_i$ чётные ($-14,\dots,14$), всего $2^{14}=16384$.
Условие точного метода $P(\SN_i\ge a)\le0.00135$ впервые выполняется
при $a=12$ (для $a=10$ имеем $P=0.006470>0.00135$, для $a=14$ $P$ равняется
$0.000061$), поэтому $a_{\mathrm{exact}}=12$,
$\mathrm{FAR}=0.001831$, $\mathrm{ARL}\approx546$. Нормальная аппроксимация:
$z_{1-\mathrm{FAR}_0/2}\sqrt{14}\approx11.22$, ближайшее достижимое чётное
значение $a_{\mathrm{norm}}=12$ --- при этом $n$ совпадает с точным.
Метод «ближайший»:
минимум $|P(|\SN_i|\ge a)-0.05|$ достигается сразу при $a=7$ и $a=8$ (оба
дают $P=0.057373$), цикл берёт меньшее, $a_{\mathrm{near}}=7$; это значение
недостижимо (значения $\SN_i$ чётные), и фактическая граница равна $8$
(рисунок~\ref{fig:modesnear}).

\begin{figure}[htbp]\centering
\begin{tikzpicture}
\begin{axis}[myaxis,
  xlabel={$a$}, ylabel={$P(|\SN|\ge a)$},
  xmin=0, xmax=15, ymin=0, ymax=0.85,
  legend pos=north east,
  title={Критерий «ближайший» при $n=14$}]
\addplot+[ycomb,mark=*,mark size=1.8pt,color=accent] coordinates {\verNearAbs14};
\addlegendentry{$P(|\SN|\ge a)$}
\addplot[dashed,alert,line width=1pt] coordinates {(0,0.05)(15,0.05)};
\addlegendentry{$\alpha=0.05$}
\addplot[only marks,mark=*,mark size=3.4pt,color=black] coordinates {\verNearMin14};
\addlegendentry{минимум $a=7,8$}
\end{axis}
\end{tikzpicture}
\caption{Критерий «ближайший» при $n=14$: двусторонний хвост
$P(|\SN|\ge a)$ и уровень $\alpha=0.05$. Кривая подходит к $\alpha$ ближе
всего при $a=7$ и $a=8$ (фактическая граница равна $8$).}
\label{fig:modesnear}
\end{figure}

Реализация трёх методов (идентификаторы совпадают с исходником):

\begin{lstlisting}[language=Python]
# 1. Conservative: smallest attainable parity value with P(SN>=a) <= FAR0/2
@classmethod
def _calc_a_exact(cls, n, far0):
    target = far0 / 2
    start = 1 if n % 2 == 1 else 2
    for a in range(start, n + 1, 2):
        if cls._p_sn_ge_a(n, a) <= target:
            return a, 2.0 * cls._p_sn_ge_a(n, a)
    return n, 2.0 * cls._p_sn_ge_a(n, n)     # fallback: widest limits

# 2. Nearest to alpha (no parity constraint -> may be non-attainable)
@classmethod
def _calc_a_nearest(cls, n, alpha):
    best_a, best_diff = 1, float('inf')
    for a in range(1, n + 1):
        diff = abs(cls._p_abs_sn_ge_a(n, a) - alpha)
        if diff < best_diff:
            best_diff, best_a = diff, a
    return best_a, cls._p_abs_sn_ge_a(n, best_a)

# 3. Symmetric: normal approximation, smallest attainable a >= z * sqrt(n)
@classmethod
def _calc_a_normal(cls, n, far0):
    z = stats.norm.ppf(1.0 - far0 / 2.0)
    a0 = int(np.ceil(z * np.sqrt(n)))
    a = a0 if a0 % 2 == n % 2 else a0 + 1
    a = min(max(a, 1), n)
    return a, 2.0 * cls._p_sn_ge_a(n, a)
\end{lstlisting}

% ================================================================== %
\section{Программная реализация}
% ================================================================== %

\subsection{Общая структура и зависимости}

Пакет \texttt{nonparametric-charts} доступен в
репозитории \path{https://github.com/nebolsinvasili/nonparametric-charts.git}. 
Основной код реализации расположен в каталоге \texttt{src/}:
\begin{center}
\footnotesize
\begin{tabular}{@{}p{5.4cm} p{8.2cm}@{}}
\toprule
Файл & Назначение \\
\midrule
\path{base/base_chart.py}            & \texttt{BaseControlChart}: абстрактный базовый класс \\
\path{charts/sign_chart.py}          & \texttt{SignChart}: знаковая карта (предмет анализа) \\
\path{charts/signed_rank_chart.py}   & \texttt{SignedRankChart}: знаково-ранговая карта \\
\path{utils/data_loader.py}          & загрузка данных процесса \\
\path{utils/data_preprocessor.py}    & подготовка данных \\
\path{utils/plot_style.py}           & единый стиль графиков \\
\bottomrule
\end{tabular}
\end{center}
Внешние зависимости: \texttt{pandas} (табличные данные),
\texttt{numpy} (векторные вычисления), \texttt{scipy}
(биномиальные распределения и специальные функции),
\texttt{matplotlib} (графическое представление). Набор зависимостей зафиксирован в
\texttt{pyproject.toml}.

\subsection{Базовый класс \texttt{BaseControlChart}}

Базовый класс (файл \path{base/base_chart.py}) задаёт единый набор методов:
абстрактные \texttt{calculate\_limits}, \texttt{plot}, \texttt{get\_summary},
а также конкретные \texttt{get\_data}, \texttt{get\_control\_limits},
\texttt{get\_violations}. В конструкторе загружаются коэффициенты контрольных
карт из \texttt{coefficients.json} (метод \texttt{\_load\_coefficients};
поиск по нескольким путям, при необходимости --- через переменную
окружения) и настраивается стиль
графиков (\texttt{\_setup\_style} через \path{utils/plot_style.py}). Здесь же
реализованы проверка правил Вестерн Электрик
(\texttt{check\_western\_electric\_rules}) и вспомогательные построения зон.
Отметим, что для непараметрических карт коэффициенты из
\texttt{coefficients.json} не используются: они относятся к параметрическим
картам (в данной работе не рассмариваются).

\subsection{Класс \texttt{SignChart}: расчёт статистик по подгруппам}

Конструктор принимает таблицу данных, объём подгруппы $n$ и целевое значение
\texttt{target\_value}, а также метод выбора константы $a$
(\texttt{approach}, \texttt{far0}, \texttt{alpha}) и необязательное явное
значение границы \texttt{a\_value}. Если целевое значение не задано, оно оценивается
как медиана всех наблюдений (метод \texttt{\_calculate\_median}); это
простейшая оценка Фазы~I.

Статистика по подгруппе вычисляется методом
\texttt{\_calculate\_signs\_for\_row}: массив значений строки очищается от
пропусков (\texttt{dropna}), после чего подсчитываются числа
\texttt{positive} ($>\!\mu_0$), \texttt{negative} ($<\!\mu_0$) и
\texttt{zero} ($=\mu_0$); результат \texttt{sn} есть разность первых двух.
Наблюдения, равные $\mu_0$, не относятся ни к положительным, ни к
отрицательным, что согласуется с непрерывной моделью \eqref{eq:sn}; однако
при значительной доле таких наблюдений эффективный объём подгруппы
уменьшается, и это обстоятельство следует иметь в виду. Знаковый критерий по
всем наблюдениям реализован отдельно в методе \texttt{calculate\_sign\_test}
(точное биномиальное $p$-значение и $z$-статистика).

\subsection{Класс \texttt{SignChart}: вероятности и константа}

Точная верхняя хвостовая вероятность вычисляется статическим методом
\texttt{\_p\_sn\_ge\_a}:
\begin{lstlisting}[language=Python]
t_threshold = int(np.ceil((n + a) / 2))
return float(np.sum([comb(n, m) * (0.5 ** n)
                     for m in range(t_threshold, n + 1)]))
\end{lstlisting}
Здесь используется соотношение $\SN_i \ge a \Leftrightarrow T_i \ge
\lceil (n+a)/2 \rceil$ и точные биномиальные коэффициенты
\texttt{scipy.special.comb}. Метод \texttt{\_p\_abs\_sn\_ge\_a} даёт
двустороннюю вероятность по симметрии, $2P(\SN_i\ge a)$, с ограничением
сверху единицей.

Три режима выбора $a$ реализованы методами
\texttt{\_calc\_a\_exact}, \texttt{\_calc\_a\_nearest} и
\texttt{\_calc\_a\_normal}. Метод
\texttt{\_calc\_a\_exact} перебирает $a$ с шагом 2, начиная с единицы
для нечётного $n$ и с двойки для чётного, и возвращает первое значение,
удовлетворяющее \eqref{eq:rule}, вместе с достигнутой FAR; при отсутствии
такового возвращается $a=n$. Метод \texttt{\_calc\_a\_nearest} перебирает все
целые $a\in[1,n]$ и минимизирует $|P(|\SN_i|\ge a)-\alpha|$. Метод
\texttt{\_calc\_a\_normal} реализует нормальную аппроксимацию: берёт
$z_{1-\mathrm{FAR}_0/2}\sqrt n$ и поднимает результат до ближайшего
достижимого значения той же чётности, что и $n$ (с ограничением $a\le n$),
как в~\eqref{eq:asym}. Выбор режима инкапсулирован в
\texttt{\_calculate\_a\_value}; сравнение режимов для текущего $n$ возвращает
\texttt{compare\_approaches}.

\subsection{Расчёт границ, нарушения и визуализация}

Метод \texttt{calculate\_limits} выполняет основной расчёт: последовательно
формирует таблицу статистик по подгруппам (цикл по строкам), определяет
константу $a$, вычисляет $\mathrm{FAR}$, $\mathrm{ARL}$, $\mathrm{MRL}$ и
заполняет словарь \texttt{control\_limits}, после чего вызывает
\texttt{\_find\_violations}. Последний отмечает подгруппы, для которых
выполнено условие \eqref{eq:signal}. Метод \texttt{plot} строит график $\SN_i$, 
контрольные границы, зоны $\pm a/3$ и $\pm 2a/3$ (где границе $a=3\sigma$), а также
информационный блок с $\mu_0$, $a$, FAR, ARL и MRL. Дополнительно реализован
\texttt{plot\_comparison} --- сопоставление трёх режимов по диапазону $n$.

\subsection{Фазы I и II и вычислительная сложность}

Целевое значение может быть задано как известное (случай K) либо оценено по
данным (метод \texttt{\_calculate\_median}). Обработка выполняется пакетно,
для всего кадра данных сразу. Потокового режима (обработки
поступающих подгрупп с обновлением оценок) не предусмотрено, что отличает
реализацию от средств, ориентированных на слежение в реальном времени.

Вычислительная сложность: Вычисление \texttt{\_p\_sn\_ge\_a} имеет сложность
$O(n)$; перебор в \texttt{\_calc\_a\_exact} даёт $O(n^2)$, а в
\texttt{\_calc\_a\_nearest} --- также $O(n^2)$. В режиме \texttt{normal} 
граница вычисляется по замкнутой нормальной аппроксимации за $O(1)$, и лишь
одна проверка хвоста \texttt{\_p\_sn\_ge\_a} для достигнутого FAR стоит
$O(n)$, что даёт суммарно $O(n)$. Формирование статистик по
подгруппам выполнено циклом языка Python по строкам со сложностью
$O(mn)$, причём с заметной постоянной. Измерения (таблица~\ref{tab:time})
показывают, что основной вклад вносит именно построчный цикл, тогда как
векторное вычисление статистики $T_i$ (как \texttt{(data > target).sum(axis=1)})
снижает время на порядки. Расчёт константы вне цикла по подгруппам для
$n=2..100$ занимает около $0.24$~с (точный режим), $0.03$~с
(режим \texttt{normal}) и $0.74$~с (режим \texttt{nearest}).

\begin{table}[htbp]\centering
\caption{Измерения времени работы реализации.}
\label{tab:time}
\small
\begin{tabular}{@{}l r@{}}
\toprule
Операция & Время, с \\
\midrule
\texttt{calculate\_limits}, $m=1000$, $n=16$          & 0.19 \\
\texttt{calculate\_limits}, $m=10000$, $n=16$         & 1.89 \\
векторная векторный расчёт $SN_i$, $m=10^6$ подгрупп  & 0.04 \\
константа $a$, $n=2..100$ (режим \texttt{exact})    & 0.24    \\
константа $a$, $n=2..100$ (режим \texttt{normal})     & 0.03 \\
константа $a$, $n=2..100$ (режим \texttt{nearest})    & 0.74 \\
\bottomrule
\end{tabular}
\end{table}

% ================================================================== %
\section{Реализации на языке R}
% ================================================================== %

Помимо рассматриваемой библиотеки, знаковая карта может быть реализована
независимо. Ниже рассмотрены две опубликованные библиотеки (пакеты
\texttt{npSPC} и \texttt{qcr}) и реализации, которые пользователи пишут
самостоятельно средствами базового R.

\subsection{Пакет npSPC}

Пакет \texttt{npSPC} представлен в работе \cite{Nascimento2026} и доступен в
репозитории \path{github.com/CER-UFBA/npSPC}. Согласно статье, пакет
содержит двенадцать функций для знаковых и знаково-ранговых карт, шесть из
которых относятся к случаю известных параметров (случай K), а шесть --- к
случаю неизвестных (случай U); карты повторяют построения
\cite{Chakraborti2019}, значением FAR по умолчанию принято 0.0027. Знаковая
карта реализована функцией \texttt{shewhart\_sn} (аргументы
\texttt{med = 0}, \texttt{far = 0.0027}). Статья также приводит сведения о
вычислительной стоимости: обработка $5\cdot10^6$ наблюдений для карты
Shewhart-SN занимает около $32$~с, что в среднем на 51.6\,\% меньше времени
выполнения параметрических алгоритмов пакета \texttt{qcc}.

Ядро правила границ (файл \texttt{R/shewhart\_sn.R}) имеет вид:
\begin{lstlisting}[language=R]
n_plus <- colSums(X > med); n_less <- colSums(X < med)
SN <- n_plus - n_less
quantil <- max(which(pbinom(0:nrow(X), nrow(X), 1/2,
                            lower.tail = FALSE) >= far))
quantil <- min(nrow(X), quantil + 1)
far <- 2 * pbinom(0:nrow(X), nrow(X), 1/2,
                  lower.tail = FALSE)[quantil]
ucl <- 2 * quantil - nrow(X); lcl <- -ucl
\end{lstlisting}
Существенные особенности:
\begin{itemize}[itemsep=1pt]
  \item матрица данных задаётся так, что строка есть наблюдение внутри
        подгруппы, а столбец --- подгруппа; статистика вычисляется векторно
        (\texttt{colSums});
  \item границы строятся через биномиальную статистику $T_i$ и потому всегда
        имеют ту же чётность, что и $n$;
  \item наблюдения, равные \texttt{med}, не учитываются; при их наличии
        выводится предупреждение;
  \item параметр \texttt{far} применяется как \emph{односторонний} порог:
        разыскивается наибольшее $q$ с $P(T_i\ge q)\ge\texttt{far}$, после
        чего полагается предел по $q+1$. Как следствие, при значении по
        умолчанию \texttt{far}=0.0027 двусторонняя FAR может достигать
        $2\,\texttt{far}$, то есть около $0.0054$ (раздел~\ref{sec:disc}).
\end{itemize}

\subsection{Пакет qcr}

Пакет \texttt{qcr} \cite{Flores2021} относится к управлению качеством, однако
классической знаковой карты в нём нет. Всего в пакете 47 функций; его
непараметрические средства \texttt{npqcs}, \texttt{npqcs.add},
\texttt{npqcs.Q}, \texttt{npqcs.r}, \texttt{npqcs.S} основаны на понятии
глубины данных \cite{Liu1995} и предназначены для многомерных наблюдений
(объект \texttt{npqcd}). Проверочный вызов \texttt{npqcs.r(method="Liu")} на
двумерных данных выполняется успешно, однако к одномерной знаковой карте
прямого отношения не имеет; ближайший одномерный инструмент
\texttt{qcs.np} есть карта числа несоответствий, использующая биномиальную
модель, но не знаковую статистику. Таким образом, \texttt{qcr} не выступает
конкурентом для реализации автора и представляет интерес лишь как
источник родственных (ранговых и глубинных) методов.

\subsection{Реализации средствами базового R}

Речь идёт не об опубликованных пакетах, а о скриптах, которые пишут
пользователи самостоятельно, опираясь на базовые функции R. Такие скрипты,
как правило, строятся из функций
\texttt{sign()} и \texttt{rank()}, а границы задаются по нормальной
аппроксимации. Поскольку $\operatorname{Var}(\SN_i)=n$, границы полагают
равными $\pm z_{1-\alpha/2}\sqrt{n}$. Пример:
\begin{lstlisting}[language=R]
SN <- colSums(sign(X - med))
a  <- ceiling(qnorm(1 - far / 2) * sqrt(n))   # normal approximation
signal <- which(abs(SN) >= a)
\end{lstlisting}
Расчёт (файл \path{experiments/sn/custom_sign_chart.R}) показывает, что
такое приближение систематически смещает границы: при $n=16$ получается
$a=12$ вместо результата точного метода $a=14$, а достигнутая FAR равна приблизительно
$0.0038$ вместо $0.0005$. При некоторых $n$ приближение даёт значение чужой
чётности, не являющееся достижимым для $\SN_i$.

% ================================================================== %
\section{Сопоставление реализаций}
\label{sec:cmp}
% ================================================================== %

Сопоставление по критериям приведено в таблице~\ref{tab:bench}.

\begin{table}[htbp]\centering
\caption{Сопоставление реализаций знаковой карты.}
\label{tab:bench}
\small
\begin{tabular}{@{}p{3.2cm} p{4.6cm} p{4.6cm}@{}}
\toprule
Критерий & Настоящая реализация & Пакет \texttt{npSPC} \\
\midrule
Представление данных & таблица, строка --- подгруппа; задаётся $n$ &
матрица, строка --- наблюдение, столбец --- подгруппа \\
Точность границ & точное биномиальное распределение &
точное биномиальное распределение \\
Интерпретация FAR & двусторонняя ($\mathrm{FAR}_0=0.0027$) &
односторонняя (тот же параметр на один хвост) \\
Достижимость $a$ & обеспечивается учётом чётности & обеспечивается \\
Наблюдения, равные $\mu_0$ & не учитываются (подсчитываются отдельно) &
не учитываются, выводится предупреждение \\
Графическое представление & \texttt{matplotlib}, зоны $\pm a/3$, $\pm 2a/3$ &
базовая функция \texttt{plot} \\
Режим работы & пакетная обработка кадра данных & пакетная обработка данных \\
Возвращаемый результат & границы, FAR, ARL, MRL, список нарушений &
строка со значением ARL \\
\bottomrule
\end{tabular}
\end{table}

По времени обработки пакет \texttt{npSPC}
обрабатывает $5\cdot10^6$ наблюдений приблизительно за $32$~с
\cite{Nascimento2026}, тогда как реализация автора обрабатывает
$1.6\cdot10^5$ наблюдений ($m=10^4$, $n=16$) за $1.9$~с
(таблица~\ref{tab:time}). При этом статистика реализована построчным циклом;
её векторное вычисление, как показано выше, ускоряет расчёт на порядок.

% ================================================================== %
\section{Вычислительный эксперимент}
% ================================================================== %

Эксперимент проведён в среде языка R 4.6.1. Пакет
\texttt{npSPC} собран из исходников репозитория
\cite{Nascimento2026}. Пакет \texttt{qcr} собран из исходников главного
хранилища CRAN.

\subsection{Проверка на управляемом процессе}

Сформирована выборка объёма $m=3000$ подгрупп по $n=16$ наблюдений, первые
2000 подгрупп относятся к управляемому процессу, начиная с 2001-й введён
сдвиг медианы на $1$ SD. Данные подготовлены в R
(\path{experiments/sn/cross_validation.R}) и переданы в Python
(\path{experiments/sn/cross_check.py}) через файл, что обеспечивает работу
обеих реализаций на идентичных числах.

Результаты:
\begin{itemize}[itemsep=1pt]
  \item значения $\SN_i$ совпали полностью для всех $3000$ подгрупп;
  \item при одной и той же границе множества сигналов совпали полностью;
  \item до момента сдвига пакет \texttt{npSPC} (\texttt{far}=0.0027) отметил
        14 ложных сигналов на 2000 подгрупп (FAR $\approx 0.0070$), 
        тогда как граница в режиме \texttt{exact} ($a=14$) --- 
        один сигнал (FAR $\approx 0.0005$); расхождение
        объясняется различием в интерпретации порога
        (раздел~\ref{sec:disc}).
\end{itemize}

\subsection{Мощность при сдвиге}

Пусть медиана процесса сдвинута на $\delta$ SD. Тогда вероятность того, что
отдельное наблюдение превысит $\mu_0$, равна $p=\Phi(\delta)$, где $\Phi$ ---
функция распределения стандартного нормального закона. Наблюдения
независимы, поэтому число превышений $T_i$ имеет \emph{точное} биномиальное
распределение $T_i\sim\mathrm{Bin}(n,p)$, а $\SN_i=2T_i-n$. Вероятность
сигнала (мощность) равна сумме двух биномиальных хвостов:
\begin{equation}
  P\bigl(|\SN_i|\ge a\bigr)=P(T_i\ge u)+P(T_i\le l),
  \qquad u=\bigl\lceil (n+a)/2\bigr\rceil,\quad
  l=\bigl\lfloor (n-a)/2\bigr\rfloor .
  \label{eq:power}
\end{equation}
Это замкнутая формула: она не использует случайных чисел и не опирается на
асимптотическое приближение. Её вычисление --- в файле
\path{experiments/sn/sn_analysis.py}.

\paragraph{Проверка методом Монте-Карло.}
Независимо моделируются $N=500\,000$ подгрупп со сдвигом $\delta$; для
каждой вычисляется $\SN_i$ и берётся доля сигналов $|\SN_i|\ge a$. По
закону больших чисел эта доля сходится к аналитической вероятности; её
стандартная ошибка не превосходит $\sqrt{0.25/N}\approx0.0007$. Совпадение
имитационной оценки с формулой (таблица~\ref{tab:montecarlo}) служит
контролем корректности и формулы, и реализации.

\begin{table}[htbp]\centering
\caption{Проверка аналитической вероятности сигнала методом Монте-Карло
($n=16$, $N=5\cdot10^5$).}
\label{tab:montecarlo}
\small
\begin{tabular}{@{}l r r r r@{}}
\toprule
Режим & $a$ & $P$, Монте-Карло & $P$, аналитически & ARL, формула \\
\midrule
управляемый                            & 14 & 0.000542 & 0.000519 & 1927.5 \\
сдвиг $1$ SD                           & 14 & 0.253904 & 0.253224 &    3.95 \\
сдвиг $1$ SD (граница \texttt{npSPC})  & 12 & 0.521900 & 0.522208 &    1.92 \\
\bottomrule
\end{tabular}
\end{table}

При $n=16$ границы различаются ($a=14$ и $a=12$), что удобно для
сопоставления (таблица~\ref{tab:oc}, рисунок~\ref{fig:oc}).

\begin{table}[htbp]\centering
\caption{Вероятность сигнала и ARL при сдвиге медианы ($n=16$).}
\label{tab:oc}
\begin{tabular}{@{}r r r r r r@{}}
\toprule
$\delta$ & $\Phi(\delta)$ & $P$, $a=14$ & ARL, $a=14$
& $P$, $a=12$ & ARL, $a=12$ \\
\midrule
0.0 & 0.5000 & 0.000519 & 1927.5 & 0.004181 & 239.2 \\
0.5 & 0.6915 & 0.022228 &   45.0 & 0.087478 &  11.4 \\
1.0 & 0.8413 & 0.253224 &    3.9 & 0.522208 &   1.9 \\
1.5 & 0.9332 & 0.709673 &    1.4 & 0.913109 &   1.1 \\
2.0 & 0.9772 & 0.949718 &    1.1 & 0.994720 &   1.0 \\
\bottomrule
\end{tabular}
\end{table}

\begin{figure}[htbp]\centering
\begin{tikzpicture}
\begin{axis}[myaxis,
  xlabel={$\delta$, SD}, ylabel={$P(|\SN|\ge a)$},
  xmin=0, xmax=2, ymin=0, ymax=1,
  legend pos=south east,
  title={Вероятность сигнала при сдвиге ($n=16$)}]
\addplot+[mark=*,mark size=1.6pt,color=accent] coordinates {\verOCBook};
\addlegendentry{граница $a=14$ (правило \eqref{eq:rule})}
\addplot+[mark=triangle,mark size=1.8pt,color=alert] coordinates {\verOCNp};
\addlegendentry{\texttt{npSPC}: $a=12$ (\texttt{far}=0.0027)}
\end{axis}
\end{tikzpicture}
\caption{Зависимость вероятности сигнала от величины сдвига. Более узкая
граница пакета \texttt{npSPC} быстрее обнаруживает сдвиг, однако ценой
повышенной FAR в управляемом состоянии.}
\label{fig:oc}
\end{figure}

\subsection{Анализ расхождений}
\label{sec:disc}

В §5.1 границы при $n=16$ разошлись: при \texttt{far}=0.0027 пакет
\texttt{npSPC} дал $a=12$, а правило \eqref{eq:rule} --- $a=14$. У
расхождения две независимые причины: разная трактовка параметра
\texttt{far} и дискретность точной границы. При $n=16$ обе дают $a=12$,
поэтому их легко спутать.

\paragraph{Трактовка параметра \texttt{far}.}
Пакет \texttt{npSPC} задаёт \texttt{far} как вероятность \emph{одного}
хвоста, тогда как правило \eqref{eq:rule} нормирует \emph{двустороннюю}
FAR: требуемый односторонний хвост не превосходит $\mathrm{FAR}_0/2$.
Иначе говоря, значению пакета \texttt{far} отвечает наша номинальная
$\mathrm{FAR}_0=2\,\texttt{far}$. Поэтому при
\texttt{far}=0.00135=0.0027/2 границы пакета совпадают с правилом
\eqref{eq:rule} для всех проверенных $n=1..20$, а при значении по умолчанию
они различаются лишь при $n=13,16,19$, где граница пакета ниже, а достигнутая
им двусторонняя FAR превышает $\mathrm{FAR}_0$ (таблица~\ref{tab:limits});
у нашего правила (\texttt{exact}) достигнутая FAR всегда не больше номинала.

\begin{table}[htbp]\centering
\caption{Границы $a$ и достигнутые FAR: правило \eqref{eq:rule} (\texttt{exact})
и \texttt{npSPC} (\texttt{far}=0.0027); режим \texttt{nearest} --- при
$\alpha=0.05$ (выборочные $n$).}
\label{tab:limits}
\small
\begin{tabular}{@{}r r r r r r r@{}}
\toprule
$n$ & $a_{\mathrm{exact}}$ & $\mathrm{FAR}_{\mathrm{exact}}$
    & $a_{\mathrm{npSPC}}$ & $\mathrm{FAR}_{\mathrm{npSPC}}$
    & $a_{\mathrm{near}}$ & $\mathrm{FAR}_{\mathrm{near}}$ \\
\midrule
 8 &  8 & 0.007812 &  8 & 0.007812 & 5 & 0.070312 \\
12 & 12 & 0.000488 & 12 & 0.000488 & 7 & 0.038574 \\
13 & 13 & 0.000244 & \textbf{11} & \textbf{0.003418} & 8 & 0.022461 \\
14 & 12 & 0.001831 & 12 & 0.001831 & 7 & 0.057373 \\
15 & 13 & 0.000977 & 13 & 0.000977 & 8 & 0.035156 \\
16 & 14 & 0.000519 & \textbf{12} & \textbf{0.004181} & 7 & 0.076813 \\
17 & 13 & 0.002350 & 13 & 0.002350 & 8 & 0.049042 \\
18 & 14 & 0.001312 & 14 & 0.001312 & 9 & 0.030884 \\
19 & 15 & 0.000729 & \textbf{13} & \textbf{0.004425} & 8 & 0.063568 \\
20 & 14 & 0.002577 & 14 & 0.002577 & 9 & 0.041389 \\
\bottomrule
\end{tabular}
\end{table}

\begin{figure}[htbp]\centering
\begin{tikzpicture}
\begin{axis}[myaxis,
  xlabel={$n$}, ylabel={константа $a$},
  xmin=0, xmax=21, ymin=0, ymax=16,
  legend pos=north west,
  title={Границы по правилам}]
\addplot+[mark=*,mark size=1.4pt,color=accent] coordinates {\verACons};
\addlegendentry{правило \eqref{eq:rule}}
\addplot+[mark=square*,mark size=1.4pt,color=alert] coordinates {\verANp};
\addlegendentry{\texttt{npSPC}, \texttt{far}=0.0027}
\addplot+[mark=triangle,mark size=1.6pt,color=third] coordinates {\verANear};
\addlegendentry{режим \texttt{nearest}}
\end{axis}
\end{tikzpicture}
\caption{Зависимость константы $a$ от объёма подгруппы. Границы пакета
\texttt{npSPC} совпадают с правилом \eqref{eq:rule}, кроме $n=13,16,19$, где
они ниже; границы режима \texttt{nearest} систематически ниже обеих.}
\label{fig:limits}
\end{figure}

\paragraph{Дискретность точной границы.}
Вторая причина --- дискретность. Граница правила \eqref{eq:rule} принимает
лишь достижимые значения $\SN_i$ нужной чётности, тогда как непрерывная
нормальная аппроксимация $\lceil z\sqrt{n}\rceil$, где
$z=z_{1-\mathrm{FAR}_0/2}$, их не учитывает. При $n=16$ она даёт $a=12$ ---
на две единицы меньше точного $a=14$, а при чётном $n$ иногда попадает в
значение чужой чётности (таблица~\ref{tab:norm}). Именно эту асимптотическую
границу с подъёмом до ближайшего достижимого значения той же чётности
воспроизводит режим \texttt{normal} библиотеки, поэтому достигнутый им FAR
не обязан быть не больше $\mathrm{FAR}_0$.

\begin{table}[htbp]\centering
\caption{Границы и достигнутые FAR: правило \eqref{eq:rule} и нормальная аппроксимация.}
\label{tab:norm}
\begin{tabular}{@{}r r r r@{}}
\toprule
$n$ & $a_{\mathrm{exact}}$ & $a_{\mathrm{norm}}$ & $\delta = a_{\mathrm{norm}} - a_{\mathrm{exact}}$ \\
\midrule
10  & 10 & 10 & 0 \\
16  & 14 & 12 & $-2$ \\
20  & 14 & 14 & 0 \\
25  & 17 & 15 & $-2$ \\
30  & 18 & 17 & $-1$ \\
40  & 20 & 19 & $-1$ \\
50  & 22 & 22 & 0 \\
100 & 32 & 30 & $-2$ \\
\bottomrule
\end{tabular}
\end{table}

\subsection{Сводный эксперимент по методам}
\label{subsec:modes-exp}

Отдельный вычислительный эксперимент (файл \path{modes_experiment.py};
порядок воспроизведения --- в \path{05_modes_experiment.txt}) сопоставляет
реализацию \texttt{npSPC} и реализацию автора на одних и тех же данных: $m=3000$
подгрупп, первые $2000$ относятся к управляемому процессу, с $2001$-й
вводится сдвиг медианы на $1$ SD.

Результаты для $n=16$ приведены в таблице~\ref{tab:modes-exp}. 
Метод \texttt{exact} ($a=14$) и реализация \texttt{npSPC} при \texttt{far}=0.00135
дают те же границы; Метод нормальной аппроксимации ($a=12$) 
совпадает с реализацией на R при значении по умолчанию
\texttt{far}=0.0027. Более узкие границы ($a=12$, $FAR=0.0042$) подвержены большему числу
ложных сигналов в IC режиме, но сокращает задержку обнаружения сдвига.

\begin{table}[htbp]\centering
\caption{Сводный эксперимент по методам, $n=16$ ($m=3000$, сдвиг с $2001$-й подгруппы); 
число ложных сигналов указано для управляемого участка ($IC$) из $2000$ подгрупп.}
\label{tab:modes-exp}
\small
\begin{tabular}{@{}l r r r r@{}}
\toprule
Режим & $a$ & $\mathrm{FAR}$ & Ложные сигналы & Задержка \\
\midrule
метод \texttt{exact}                    & 14 & 0.000519 &   2 & 4 \\
метод \texttt{normal}                   & 12 & 0.004181 &   6 & 1 \\
метод \texttt{nearest} ($\alpha=0.05$)  &  7 & 0.076813 & 136 & 1 \\
метод \texttt{normal} (без чётности)    & 12 & 0.004181 &   6 & 1 \\
\texttt{npSPC} (\texttt{far}=0.0027)    & 12 & 0.004181 &   6 & 1 \\
\texttt{npSPC} (\texttt{far}=0.00135)   & 14 & 0.000519 &   2 & 4 \\
\bottomrule
\end{tabular}
\end{table}

Границы, полученные всеми режимами для трёх объёмов подгруппы, сведены в
таблице~\ref{tab:modes-alln}.

\begin{table}[htbp]\centering
\caption{Константа $a$ по режимам для $n=13,16,19$.}
\label{tab:modes-alln}
\small
\begin{tabular}{@{}l r r r@{}}
\toprule
Режим & $n=13$ & $n=16$ & $n=19$ \\
\midrule
метод \texttt{exact}                        & 13 & 14 & 15 \\
метод \texttt{normal}                       & 11 & 12 & 15 \\
метод \texttt{nearest} ($\alpha=0.05$)      &  8 &  7 &  8 \\
метод \texttt{normal} (без чётности)        & 11 & 12 & 14 \\
\texttt{npSPC} (\texttt{far}=0.0027)        & 11 & 12 & 13 \\
\texttt{npSPC} (\texttt{far}=0.00135)       & 13 & 14 & 15 \\
\bottomrule
\end{tabular}
\end{table}

% \paragraph{Что лучше.} Ответ зависит от постановки задачи. Если задан целевой
% $\mathrm{FAR}_0=0.0027$, верна наша точная граница: она обеспечивает
% $\mathrm{FAR}\le\mathrm{FAR}_0$, тогда как значение по умолчанию в
% \texttt{npSPC} из-за односторонней трактовки параметра \texttt{far}
% фактически удваивает $\mathrm{FAR}$ (до $\approx0.0042$ при $n=16$) и даёт
% втрое больше ложных тревог. Выигрыш более узкой границы --- обнаружение
% сдвига на две--три подгруппы раньше, что не окупает потери управляемости.
% При \texttt{far}=0.00135 обе реализации совпадают, так что расхождение
% связано не с качеством кода, а с интерпретацией параметра. Наш метод
% нормальной аппроксимации воспроизводит как раз асимптотическую
% границу \texttt{npSPC} по умолчанию --- это удобный режим для быстрой оценки
% и сопоставления, но не для строгого контроля, так как достигнутый FAR может
% превышать $\mathrm{FAR}_0$. Режим \texttt{nearest} в статистическом
% управлении процессами непригоден: его $\mathrm{FAR}\approx0.05$--$0.08$ даёт
% десятки ложных тревог на управляемом участке.

% ================================================================== %
\section{Результаты и рекомендации}
% ================================================================== %

\subsection{Преимущества рассматриваемой реализации}

Границы вычисляются по точному биномиальному распределению, что подтверждено
как точным расчётом, так и экспериментом Монте-Карло
(таблицы~\ref{tab:montecarlo},~\ref{tab:norm}). Реализация возвращает не
только границы, но и $\mathrm{FAR}$, $\mathrm{ARL}$, $\mathrm{MRL}$ и
перечень нарушений, тогда как \texttt{shewhart\_sn} ограничивается значением
ARL. А для работы не требуются тяжёлые внешние зависимости.

\subsection{Направления развития}

Основное направление связано с оцениванием целевого значения: если $\mu_0$
оценивается по тем же данным, распределение $\SN_i$ перестаёт быть строго
биномиальным, вследствие чего номинальные $\mathrm{FAR}$ и $\mathrm{ARL}$
могут не достигаться, и целесообразен переход к случаю U с поправкой на
оценивание \cite{Chakraborti2019}. Далее, реализацию можно дополнить
потоковым режимом Фазы II с обновлением оценок по мере поступления подгрупп
и векторизовать построение статистик, устранив построчный цикл языка Python
(таблица~\ref{tab:time}). Результат метода \texttt{nearest} следует либо
привести к достижимым значениям $a$, либо явно помечать как условный, а
обработку наблюдений, равных $\mu_0$, --- документировать и оценивать их
долю. Наконец, у пакета \texttt{npSPC} полезно заимствовать поддержку
односторонних карт, а у пакета \texttt{qcr} ---
ранговые и глубинные карты для многомерного случая.

% ================================================================== %
\section{Заключение}
% ================================================================== %

Для случая известного целевого значения реализация математически корректна:
знаковая статистика, правило границ и характеристики длины серии
соответствуют первоисточникам \cite{Amin1995, Gibbons2010}, а перекрёстный
вычислительный эксперимент с пакетом \texttt{npSPC} на идентичных данных дал
полное совпадение значений $\SN_i$ и меток сигналов. Единственное
содержательное расхождение по границам с пакетом \texttt{npSPC} объясняется
односторонней интерпретацией параметра \texttt{far}: при значении
\texttt{far}=0.0027 границы пакета оказываются уже, а FAR --- выше
номинальной; при \texttt{far}=0.00135 результаты совпадают. Классической
знаковой карты в пакете \texttt{qcr} нет, вследствие чего по данному виду
карт рассматриваемая реализация не имеет прямых аналогов среди разобранных
средств, а по точности превосходит типичные реализации средствами базового
R с нормальной аппроксимацией. Основным направлением совершенствования
является векторизация вычислений и учёт эффекта оценивания целевого значения.

% ================================================================== %
\section{Приложения}
% ================================================================== %

\paragraph{Сценарии эксперимента} (каталог \path{experiments/sn/}):
\begin{itemize}[itemsep=1pt]
  \item \path{sn_analysis.py} --- точные вероятности, проверка Монте-Карло,
        измерения времени;
  \item \path{cross_validation.R} и \path{cross_check.py} --- перекрёстный
        эксперимент на идентичных данных;
  \item \path{modes_experiment.py} и \path{mode_comparison.R} --- сводный
        эксперимент по режимам с вызовом пакета \texttt{npSPC};
  \item \path{custom_sign_chart.R} --- реализация средствами базового R и
        нормальная аппроксимация;
  \item \path{qcr_check.R} --- проверка средств пакета \texttt{qcr};
  \item \path{gen_figdata.py} --- координаты графиков;
  \item \path{00_ENVIRONMENT.txt} и \path{0X_*.txt} --- инструкции по
        воспроизведению каждого эксперимента.
\end{itemize}

\paragraph{Воспроизведение среды языка R} (без прав суперпользователя):
\begin{lstlisting}
micromamba create -p ./renv -c conda-forge r-base gfortran
micromamba install -p ./renv -c conda-forge r-dplyr r-purrr r-fda r-ks \
    r-nlme r-mgcv r-lattice r-survival r-mvtnorm r-mass r-matrix \
    r-foreach r-iterators r-doparallel r-knitr r-ksamples r-suppdists \
    r-goftest
R CMD INSTALL npSPC-main        # npSPC package (CER-UFBA/npSPC)
curl -LO https://cloud.r-project.org/src/contrib/qcc_2.7.tar.gz
curl -LO https://cloud.r-project.org/src/contrib/fda.usc_2.2.0.tar.gz
curl -LO https://cloud.r-project.org/src/contrib/qcr_1.4.tar.gz
R CMD INSTALL qcc fda.usc qcr   # in this order
Rscript experiments/sn/cross_validation.R
\end{lstlisting}

\paragraph{Листинг расчёта константы $a$:}
\begin{lstlisting}[language=Python]
@classmethod
def _calc_a_exact(cls, n, far0):
    target = far0 / 2
    start = 1 if n % 2 == 1 else 2
    for a in range(start, n + 1, 2):
        if cls._p_sn_ge_a(n, a) <= target:
            return a, 2.0 * cls._p_sn_ge_a(n, a)
    return n, 2.0 * cls._p_sn_ge_a(n, n)   # fallback: widest limits
\end{lstlisting}

% ================================================================== %
\begin{thebibliography}{9}
% ================================================================== %

\bibitem{Amin1995}
Amin, R.~W., Reynolds~Jr., M.~R., Bakir, S.~T. (1995).
Nonparametric quality control charts based on the sign statistic.
\emph{Communications in Statistics: Theory and Methods}, 24(6), 1597--1623.

\bibitem{Bakir2004}
Bakir, S.~T. (2004).
A distribution-free Shewhart quality control chart based on signed-ranks.
\emph{Quality Engineering}, 16(4), 613--623.

\bibitem{Bakir1979}
Bakir, S.~T., Reynolds~Jr., M.~R. (1979).
A nonparametric procedure for process control based on within-group ranking.
\emph{Technometrics}, 21(2), 175--183.

\bibitem{Chakraborti2019}
Chakraborti, S., Graham, M.~A. (2019).
\emph{Nonparametric Statistical Process Control}. Wiley.

\bibitem{Gibbons2010}
Gibbons, J.~D., Chakraborti, S. (2010).
\emph{Nonparametric Statistical Inference}, 5th ed. CRC Press.

\bibitem{Nascimento2026}
Nascimento, D.~C., Fons\^eca, A.~O., Fons\^eca, M.~B., Silva, M.~A.~F.,
Ferreira, P.~H. (2026).
Nonparametric SPC: Solving complex industrial monitoring tasks with low
computational cost. \emph{Array}, 30, 100828.
\url{https://doi.org/10.1016/j.array.2026.100828}.

\bibitem{Flores2021}
Flores, M., Fern\'andez-Casal, R., Naya, S., Tarr\'io-Saavedra, J. (2021).
qcr: Quality Control Review. \emph{The R Journal}.
\url{https://doi.org/10.32614/RJ-2021-034}.

\bibitem{Liu1995}
Liu, R.~Y. (1995).
Control charts for multivariate processes.
\emph{Journal of the American Statistical Association}, 90(432), 1380--1387.

\end{thebibliography}

\end{document}
